% Appendix G, from exam/paper_b_solutions.md: the model answers to Paper B, with marks.

% The answers' headings are labelled ansb:q:N.
\renewcommand{\answerprefix}{ansb}

\chapter{Model Answers to Paper B}
\label{ansb:ch:answers}

Marks are shown per part. Method carries them: a correct derivation with a slip in it is worth more
than a correct answer with no working, and later parts consume earlier ones, so an error should be
followed through rather than penalised twice. A candidate whose register map is wrong in Question~5
and who then uses their own map consistently in Question~6 loses the marks once.

Where a part asks for VHDL, mark the \textbf{hardware described}, not the syntax: the sensitivity
list, the reset branch, whether a pulse is one clock wide, whether every path assigns its output,
and which way a shift goes. Where it asks for C++, mark the \textbf{byte order, the ordering of
register accesses, and the qualifiers that carry meaning} --- \code{const}, \code{noexcept},
\code{override}, and what has been deleted --- not the punctuation.

Where a part asks for a review, naming the defect earns half and saying what it does to the running
system earns the other half. The questions are in Appendix~\ref{exb:ch:paper}.

% ==========================================================================================
\answersection{1}{Building the top before anything it holds}{12}

\answerpart{a}{3}
\begin{vhdlcode}
--------------------------------------------------------------------------------
-- UART peripheral top level.
--
-- Inputs:
--    - clock: 50 MHz system clock.
--    - reset_n: Asynchronous active-low reset.
--    - sclk, mosi, ss: SPI from the processor.
--    - rx: UART serial input.
--
-- Outputs:
--    - miso: SPI to the processor.
--    - tx: UART serial output.
--------------------------------------------------------------------------------
library ieee;
use ieee.std_logic_1164.all;

entity uart_top is
    port(clock, reset_n: in std_logic;
         sclk, mosi, ss: in std_logic;
         rx            : in std_logic;
         miso, tx      : out std_logic);
end entity;
\end{vhdlcode}

{\itshape (2 marks: 1 for the eight ports in the right order with the right directions, 1 for the
inputs-first grouping, the clauses and the banner.)\par}

\medskip
\textbf{The clause it does not need.} \code{ieee.numeric_std}. Nothing in the skeleton of
\chapref{1} converts between a vector and an integer: the top wires bytes and vectors from block to
block and never interprets one as a number.

It is added in \textbf{\chapref{2}}, for the single guarded conversion that turns the register
bank's \code{BAUD_DIV} vector into the \code{natural} \code{baud_gen} expects:

\begin{vhdlcode}
baud_div_int <= 1 when is_x(baud_div) or unsigned(baud_div) = 0
              else to_integer(unsigned(baud_div));
\end{vhdlcode}

\noindent\code{unsigned} and \code{to_integer} come from \code{numeric_std}; \code{is_x} comes from
\code{std_logic_1164}, so it needs no clause of its own. That expression is the only place in
\code{uart_top} where a \code{std_logic_vector} becomes an integer.

\textbf{Why never \code{use work.uart_def.all}.} The top carries the register bus ---
\code{reg_addr}, \code{reg_wdata}, \code{reg_write}, \code{reg_rdata} --- from the bridge to the
bank without ever \textbf{naming} a register or a bit on it. It is wiring, not interpretation.
\code{uart_regs} is the only synthesizable module that interprets what is on that bus, so it is the
only one that needs the package (the testbenches use it too, mostly for \code{to_hex}). \worth{1}

\answerpart{b}{5}
\textbf{The signals.} Read each entity and declare, for every port that does not connect straight to
a \code{uart_top} port, a signal of the same type:

\begin{vhdlcode}
architecture behaviour of uart_top is
signal reset_s2_n   : std_logic;                      -- from reset_sync, to everything
signal spi_rx_data  : std_logic_vector(7 downto 0);   -- slave  -> bridge
signal spi_rx_valid : std_logic;                      -- slave  -> bridge
signal spi_ss_active: std_logic;                      -- slave  -> bridge
signal spi_tx_data  : std_logic_vector(7 downto 0);   -- bridge -> slave
signal reg_addr     : std_logic_vector(3 downto 0);   -- bridge -> uart_regs (L05)
signal reg_wdata    : std_logic_vector(31 downto 0);  -- bridge -> uart_regs (L05)
signal reg_write    : std_logic;                      -- bridge -> uart_regs (L05)
signal reg_rdata    : std_logic_vector(31 downto 0);  -- uart_regs (L05) -> bridge
begin
\end{vhdlcode}

{\itshape (2 marks: 1 for the nine signals and their types, 1 for the \code{spi_} prefix and
\code{_s2} suffix used correctly and for declaring nothing a later chapter brings.)\par}

\medskip
\textbf{The instantiations.} \textit{(2 marks)}

\begin{vhdlcode}
    reset_sync: entity work.reset_sync
        port map(clock, reset_n, reset_s2_n);

    spi_slave: entity work.spi_slave
        port map(clock, reset_s2_n, sclk, mosi, ss, spi_tx_data,
                 miso, spi_rx_data, spi_rx_valid, spi_ss_active);

    spi_reg_bridge: entity work.spi_reg_bridge
        port map(clock, reset_s2_n, spi_ss_active, spi_rx_data, spi_rx_valid, reg_rdata,
                 spi_tx_data, reg_addr, reg_wdata, reg_write);
\end{vhdlcode}

\noindent Note that \code{miso} connects \textbf{straight to the entity port}, so it needs no signal
of its own --- that is the reason the list above has nine entries rather than ten, and it is the
whole method of the exercise.

\textbf{The two \code{rx_data} ports.} \textit{(1 mark)} They are not one bus, and they are not even
the same direction. \code{spi_slave}'s \code{rx_data} is an \textbf{output}: bytes it has finished
shifting in from \code{MOSI}. The bridge's \code{rx_data} is an \textbf{input}: the bytes it
assembles into transactions. So \code{spi_rx_data} runs \textbf{slave to bridge}. \code{spi_tx_data}
runs the other way, from the bridge's \code{tx_data} \emph{output} into the slave's \code{tx_data}
\emph{input}, carrying the bytes to be shifted out on \code{MISO}. Read them as two one-directional
lanes that happen to share a naming convention, and the wiring falls out.

\textbf{Why \code{reg_addr} is four bits.} The protocol's command byte reserves bits 3--0 for the
register index, so the master can send any value from 0 to 15 and the bus has to be able to carry
it. Indices 7--15 are \emph{defined} as reserved --- a read returns \code{0x00000000} and a write is
ignored --- which is a rule the bank can only implement if it can see the out-of-range index in the
first place. Three bits would make an illegal address indistinguishable from a legal one.

\answerpart{c}{2}
\begin{vhdlcode}
    reg_rdata <= (others => '0');
\end{vhdlcode}

{\itshape (Half a mark --- the line is trivial; the marks are in what follows.)\par}

\medskip
\textbf{With it in place}, an SPI read transaction returns \textbf{\code{0x00000000}} for every
register index. That is harmless until \chapref{5} because there is no register bank yet: nothing in
the design has a value to return, and no testbench asks for one --- \code{uart_top_tb} is
\emph{skipped} until its datapath exists, and the two transport testbenches drive the provided
blocks directly. \worth{1}

\textbf{Without it}, \code{reg_rdata} has \textbf{no driver at all}. That is legal VHDL and the file
still analyzes. But an undriven signal of a resolved type sits at the \textbf{left-most value of its
type}, which for \code{std_logic} is \code{'U'}, \emph{uninitialized}. So ``returns'' is the wrong
word: the bridge latches all-\code{'U'}, shifts \code{'U'} out on \code{MISO}, and the metavalue
propagates into anything that looks at it. It is not zero and it is not garbage-but-defined; it is
the simulator's marker for ``nothing has ever driven this'', which is exactly the information the
placeholder exists to suppress while the bank is missing. \textit{(Half a mark.)}

The other half of the mark is for saying \textbf{when it must go}: in \chapref{5}, in the same edit
that instantiates \code{uart_regs}. Leave it and the vector has two drivers.

\answerpart{d}{2}
\textbf{What the top fixes.} Every block's \textbf{port contract}: the number of ports, their order,
their directions and their types. The moment \code{uart_top} writes
\code{port map(clock, reset_s2_n, baud_div_int, baud_tick)}, \code{baud_gen} has no interface
decisions left to make --- and neither does its testbench, which was written against the same
contract before either existed. Building top-down means a block arrives into a shape that is already
waiting for it, rather than the top being rewritten around each block as it lands. \worth{1}

\textbf{What has to be true of \code{baud_gen}.} Its entity must declare exactly four ports, in this
order and with these types:

\begin{narrowtable}{@{}rll@{}}
  \thead{\#} & \thead{Dir} & \thead{Type} \\
  \midrule
  1 & in & \code{std_logic} \\
  2 & in & \code{std_logic} \\
  3 & in & \code{natural range 1 to 65535} \\
  4 & out & \code{std_logic} \\
\end{narrowtable}

\textbf{What is not part of the contract: the names.} \file{hw/README.md} says so outright --- rename
every port and everything still binds, because nothing in this course associates a port by name. The
chapters keep their names anyway so that the prose and the code talk about the same signals, but
that is a convenience for readers, not a rule the tools enforce. It is also, as Question~\ref{exa:q:1}
of Paper~A explores, the reason a transposed pair of same-type ports is invisible until the system
testbench runs. \worth{1}

% ==========================================================================================
\answersection{2}{The generator, the divider, and the frame}{13}

\answerpart{a}{4}
\begin{vhdlcode}
--------------------------------------------------------------------------------
-- Baud rate generator.
--
-- Inputs:
--    - clock: 50 MHz system clock.
--    - reset_s2_n: Active-low synchronized reset.
--    - div: Baud rate divider. Must be greater than 0.
--
-- Outputs:
--    - tick: Baud rate tick.
--------------------------------------------------------------------------------
library ieee;
use ieee.std_logic_1164.all;

entity baud_gen is
    port(clock, reset_s2_n: in std_logic;
         div              : in natural range 1 to 65535;
         tick             : out std_logic);
end entity;

architecture behaviour of baud_gen is
signal counter: natural;
begin
    process(clock, reset_s2_n) is
    begin
        if (reset_s2_n = '0') then
            tick    <= '0';
            counter <= 0;
        elsif (rising_edge(clock)) then
            if (counter >= div - 1) then
                tick    <= '1';
                counter <= 0;
            else
                tick    <= '0';
                counter <= counter + 1;
            end if;
        end if;
    end process;
end architecture;
\end{vhdlcode}

{\itshape (3 marks: 1 for the entity with \code{div} as a \code{natural} and the four ports in order;
1 for the asynchronous reset with \code{reset_s2_n} in the sensitivity list and both outputs cleared;
1 for a \textbf{registered}, one-clock-wide \code{tick} with the counter reloading to zero.)\par}

\medskip
Deduct for a \code{tick} driven combinationally, for a synchronous reset, or for a counter that
reloads to \code{1}.

\textbf{The comparison.} \textit{(1 mark)}

At \code{div = 1}, \code{div - 1} is \code{0} and the counter is always \code{0} when the test
runs, so both forms fire on \textbf{every} clock. They are indistinguishable, and the same is true
for any stable \code{div}: the counter reloads the instant it reaches the threshold and never passes
it.

At \code{div = 0} they diverge completely. \code{div - 1} evaluates to \code{-1}:
\begin{itemize}
  \item \code{counter >= -1} is \textbf{always true}, so the block ticks every cycle and the counter
    never leaves zero. Useless output, but the design keeps running and stays inside its declared
    range.
  \item \code{counter = -1} is \textbf{never true}, so the counter is incremented on every edge for
    ever. \code{counter} is a \code{natural}, so this ends in a \textbf{run-time bound check failure}
    that aborts the simulation, and in synthesis it is a counter that free-runs and wraps.
\end{itemize}
So \code{>=} is the one that keeps the counter inside its declared range. It also corrects a
\code{div} changed mid-count on the next edge rather than after a wrap. One character of typing for
both.

Worth noting for a candidate who raises it: \code{div} is declared \code{natural range 1 to 65535},
so \code{0} cannot legally be \emph{bound} to it. The \code{>=} is defensive against the case the
subtype is supposed to rule out, which is exactly when defensiveness is cheap and worth having.

\answerpart{b}{3}
\begin{console}
exact divider = 50_000_000 / (16 x 9600) = 50_000_000 / 153_600 = 325.52
BAUD_DIV       = round(325.52) = 326

tick rate = 50_000_000 / 326 = 153_374.2 Hz
baud rate = 153_374.2 / 16   = 9585.9 baud
error     = (9585.9 - 9600) / 9600 = -0.15%   (slow)
\end{console}

{\itshape (2 marks: 1 for the divider, 1 for the achieved rate and a correctly signed error.)\par}

\medskip
\textbf{Why \code{16 * baud}.} \textit{(1 mark)} The receiver cannot see where a bit begins ---
there is no clock on the wire, only the line itself. So it \textbf{oversamples}: it watches
\code{rx} sixteen times per bit, finds the falling edge that starts a frame, and then aims for the
middle of each bit rather than its edge. For the transmitter and the receiver to share \textbf{one}
time base rather than each keeping its own, that base has to tick at the finer rate.

\textbf{\code{uart_rx} spends the sixteen.} It uses them for resolution: one to notice the start
edge, eight to reach the middle of the start bit for the re-check that rejects a glitch, and then one
sample every sixteen thereafter, at each bit's centre.

\textbf{\code{uart_tx} does not need them} and simply counts: sixteen ticks to a bit, changing the
line only on bit boundaries. It shares the generator not because it needs the resolution but so that
both halves are paced by the same divider from the same clock, with no handshake between them.

\answerpart{c}{3}
\textbf{The frame.} \textit{(1 mark)}

\begin{vhdlcode}
frame <= STOP_BIT & parity_bit & data & START_BIT;   -- frame(10 downto 0)
\end{vhdlcode}

\begin{narrowtable}{@{}lccccc@{}}
  \thead{Index} & 0 & 1 .. 8 & 9 & 10 \\
  \midrule
  Carries & start & \code{data(0)} .. \code{data(7)} & \textbf{parity} & stop \\
\end{narrowtable}

The vector is now 11 bits and \code{bit_idx} is declared \code{natural range 0 to 10}; the terminal
test compares \code{bit_idx} with the last index of the frame actually being sent, 10 with parity and
9 without. The parity bit goes \textbf{between the last data bit and the stop bit}, which is where
the line protocol puts it, and the stop bit moves from index 9 to index 10.

\textbf{The parity bit.} \textit{(1 mark)}

\begin{vhdlcode}
even_parity <= data(0) xor data(1) xor data(2) xor data(3)
           xor data(4) xor data(5) xor data(6) xor data(7);

parity_bit  <= even_parity when parity_odd = '0' else not even_parity;
\end{vhdlcode}

\noindent Even parity makes the total number of ones --- data plus parity --- \textbf{even}, and the
XOR of the data bits is exactly the bit that achieves it. Odd parity is its complement. (With
\code{parity_en} clear the parity field is not sent at all: the stop bit stays at index 9 and the
frame ends there, exactly as in 8N1. A design that always sends eleven bits and drives the parity
position \code{'0'} when parity is off is \emph{wrong}: an 8N1 receiver samples that \code{'0'} where
it expects the stop bit, and reports a framing error on every byte.)

\textbf{The variable field, and the cost.} \textit{(1 mark)} With a \code{two_stop} input, the
\textbf{stop field} is the only field whose length varies --- 1 or 2 bits. Parity is a different
matter: it adds a field rather than lengthening one.

Frame lengths: \textbf{10} bits (8N1), \textbf{11} (either parity or a second stop bit), \textbf{12}
(both).

At 115200 baud, throughput is \code{baud / bits-per-frame}:

\begin{console}
10 bits: 115_200 / 10 = 11_520 bytes/s
12 bits: 115_200 / 12 =  9_600 bytes/s
\end{console}

\noindent so the longest frame costs \textbf{1920 bytes/s}, a drop of \textbf{16.7\,\%}, for the
same line rate. The wire runs at exactly the same speed; you are simply spending a third of it on
framing (4 bits in 12) rather than data, where 8N1 spent a fifth (2 bits in 10).

\answerpart{d}{3}
\textbf{What the testbench checks about \code{busy}.} \textit{(1 mark)} Nothing. \code{uart_tx_tb}
\textbf{binds} \code{busy} and \code{done} but never asserts on either --- \secref{c2:app:b} says so
in as many words. What it does check is that \code{tx} idles high, goes low within two bit periods
of \code{start}, that the start bit is low at its centre, that the eight data bits match least
significant first, and that the stop bit is high. A \code{busy} that rises one clock late is
entirely invisible to it.

\textbf{What the feeder does.} \textit{(1 mark)} The feeder is

\begin{vhdlcode}
tx_load <= (not tx_empty) and (not tx_busy);
tx_pop  <= tx_load;
\end{vhdlcode}

\noindent On the edge that loads a byte, the FIFO is popped and the transmitter enters
\code{STATE_SEND}. On the \textbf{next} clock, \code{tx_busy} is still low --- it rises a clock late
--- and \code{tx_empty} is still low if any byte remains queued. So \code{tx_load} is high for a
\textbf{second} clock, producing a second \code{start} pulse and a second \code{tx_pop}.

\textbf{What happens to the byte.} \code{uart_tx} ignores \code{start} outside \code{STATE_IDLE}, so
the second start does nothing --- the candidate is right about that, and it is what makes the bug
survive review. But \code{tx_pop} is not ignored: the FIFO honours it and \textbf{discards its front
byte}. So the byte behind the one in flight is popped and never transmitted, once per frame.

Note the edge case that hides it: with exactly one byte queued, the FIFO is already empty after the
first pop, so \code{tx_empty} is high on the second clock and \code{tx_load} is low. The bug only
appears from the second queued byte onward.

\textbf{The fix and the first testbench that could catch it.} \textit{(1 mark)}

\begin{vhdlcode}
busy <= '1' when state = STATE_SEND else '0';
\end{vhdlcode}

\noindent A concurrent decode of the state, so \code{busy} is high in the same simulation cycle as
the state change rather than a clock later.

The first testbench that \textbf{could} have caught it is \code{uart_top_tb}, in \textbf{\chapref{5}}
--- the first testbench in which the feeder, the FIFO and the transmitter exist at the same time. Be
honest about the rest: as written it writes exactly \textbf{one} byte to \code{TX_DATA}, so it lands
squarely in the edge case above and does \textbf{not} catch it. Catching it needs a testbench that
queues two bytes behind one in flight and checks that all of them leave. Full marks for the honest
answer; half for ``the \code{uart_top_tb} of \chapref{5}'' without the caveat.

% ==========================================================================================
\answersection{3}{The asynchronous line, written out}{14}

\answerpart{a}{4}
\begin{vhdlcode}
--------------------------------------------------------------------------------
-- Generic double-flop synchronizer.
--
-- Generics:
--    - COUNT: Number of inputs to synchronize. Must be greater than 0.
--
-- Inputs:
--    - clock: 50 MHz system clock.
--    - reset_s2_n: Active-low synchronized reset.
--    - async_in: Inputs to synchronize.
--
-- Outputs:
--    - sync_out: Synchronized signals.
--------------------------------------------------------------------------------
library ieee;
use ieee.std_logic_1164.all;

entity sync is
    generic(COUNT: natural range 1 to 15 := 1);
    port(clock, reset_s2_n: in std_logic;
         async_in         : in std_logic_vector(COUNT-1 downto 0);
         sync_out         : out std_logic_vector(COUNT-1 downto 0));
end entity;

architecture behaviour of sync is
signal sync_s1, sync_s2: std_logic_vector(COUNT-1 downto 0);
begin
    sync_out <= sync_s2;

    process(clock, reset_s2_n) is
    begin
        if (reset_s2_n = '0') then
            sync_s1 <= (others => '0');
            sync_s2 <= (others => '0');
        elsif (rising_edge(clock)) then
            sync_s1 <= async_in;
            sync_s2 <= sync_s1;
        end if;
    end process;
end architecture;
\end{vhdlcode}

{\itshape (2 marks: 1 for the generic with its default and the vector ports sized from it; 1 for two
flops in series with the \textbf{asynchronous} reset and \code{sync_out} taken from the
\textbf{second}.)\par}

\medskip
The commonest error worth deducting for is \code{sync_out <= sync_s1}, which throws away the entire
point: the output must come from the flop that has had a full clock period to settle.

\textbf{Two edges.} \textit{(1 mark)} \code{sync_out} follows \code{async_in} after \textbf{exactly
two} rising edges.
\begin{itemize}
  \item \textbf{One is not a synchronizer.} A single register gives a metastable output no settling
    time at all: whatever it is doing at the next edge is what the design consumes. It looks
    identical in simulation --- nothing in a simulator is ever metastable --- which is precisely why
    \code{sync_tb}'s ``exactly two edges'' check exists and why deleting a flop is the bug that
    testbench is \emph{for}.
  \item \textbf{Three is not wrong, it is unspecified.} An extra stage buys more settling time and
    would be right at a much higher clock or for a far more hostile input, but here it adds a clock
    of latency, fails \code{sync_tb}, and changes the receiver's edge-detection timing. The contract
    is two, so the module delivers two.
\end{itemize}

\textbf{A \code{COUNT = 8} instance.} \textit{(1 mark)} Each bit crosses \textbf{independently},
through its own pair of flops. Two bits that change on the same source edge can therefore resolve on
different clocks, so \code{sync_out} may show one bit's new value beside another's old one for a
cycle.

That is \textbf{not a defect in \code{sync}}: each bit individually is correctly synchronized, which
is everything a synchronizer promises. The promise was never coherence across bits.

So the kind of signal you must never put through a wide instance is a \textbf{multi-bit value whose
bits must stay coherent} --- a counter, an address, any encoded number crossing clock domains. The
receiving side can read a value the source never held: a counter going from \code{0111} to
\code{1000} can be sampled as anything from \code{0000} to \code{1111}. Use a Gray code (only one bit
changes per step, so the worst case is the old value or the new one), a request/acknowledge
handshake, or an asynchronous FIFO. \code{rx} is a single bit, so the question does not arise here
--- but it arises the first time somebody reaches for a wide instance.

\answerpart{b}{4}
\begin{vhdlcode}
when STATE_IDLE =>
    -- Leave idle only on a high-to-low transition, never on a level.
    if ((rx_s2 = START_BIT) and (rx_prev = STOP_BIT)) then
        state <= STATE_START;
        ticks <= 0;              -- restart the count from this edge
    end if;

when STATE_START =>
    -- Half a bit on from the edge, check the line is still low.
    if (sample_start_bit = '1') then
        if (rx_s2 = '0') then
            state <= STATE_DATA;
            ticks <= 0;
        else
            state <= STATE_IDLE; -- a glitch, not a start bit
        end if;
    end if;
\end{vhdlcode}

\noindent and, \textbf{last on the tick}, after the \code{case} and inside the
\code{if baud_tick = '1'}:

\begin{vhdlcode}
rx_prev <= rx_s2;
\end{vhdlcode}

{\itshape (2 marks: 1 for the edge test and the glitch re-check with both outcomes; 1 for the history
flop and its placement.)\par}

\medskip
\textbf{Why the placement matters.} \code{rx_prev} is updated \textbf{on the tick}, inside
\code{if baud_tick = '1'}, which is what makes the comparison ``this tick against the previous
tick''. Update it on every clock instead and it holds the line as it was one clock ago, not one tick
ago: a falling edge is then seen only if it arrives in the last clock before a tick, so the receiver
misses nearly every start bit, and \code{uart_rx_tb} fails its first case, ``expected exactly one
valid byte''. Where it sits \emph{within} the tick does not change the hardware: \code{rx_prev} is a
signal, so the \code{case} reads the value it had before this clock edge whatever the textual order,
and a receiver with the assignment moved above the \code{case} passes \code{uart_rx_tb} unchanged.
Writing it last is for the reader, since it puts the update where it takes effect. (Were it a
\emph{variable}, writing it first really would compare \code{rx_s2} with itself, and the receiver
would never leave idle.) Worth noting alongside: the \code{ticks <= 0} in the idle branch overrides
the tick-counter block above it, because that block comes first in the process and VHDL gives the
\textbf{last} assignment to a signal in a process the final say.

\textbf{The level test and the break.} \textit{(2 marks)}

With \code{if rx_s2 = '0' then}, the receiver re-arms the moment the previous frame's stop bit has
been judged. During a break --- the line held low for longer than a frame, which is what a
transmitter losing power or a cable coming loose looks like --- it marches through back-to-back
\textbf{all-zero frames}. Each of those ends on a low stop bit and is correctly rejected as a framing
error, and that part is fine.

The failure is the frame still \textbf{in flight when the break ends}. Its stop bit is sampled after
the line has returned high, so as far as the receiver can tell it is a well-formed frame:
\code{valid} pulses and a byte assembled out of nothing --- the tail of the break plus the returning
idle --- is pushed into the RX FIFO as real data. Software receives a byte nobody sent.

With an edge test there is no second falling edge until the line has gone high again, so no frame is
ever in flight across that boundary. One flip-flop, and it is the difference between a receiver that
works on clean data and one that survives a cable being unplugged.

\textbf{The testbench case} is \code{uart_rx_tb}'s break case: it holds the line low for longer than
one frame and then idles long enough for a frame in flight to complete. Its message names the fix
outright --- ``a break produced a byte - leave \code{STATE_IDLE} on a falling edge'' --- and a
second check requires the break to be reported as at least one framing error, so a receiver that goes
silent instead of wrong does not pass either.

\answerpart{c}{3}
\begin{vhdlcode}
    elsif (rising_edge(clock)) then
        -- Defaults first: any pulse from the previous cycle lasts exactly one clock.
        data_out  <= (others => '0');
        valid     <= '0';
        frame_err <= '0';
        ...
            when STATE_STOP =>
                if (sample_data = '1') then
                    if (rx_s2 = '1') then
                        data_out <= frame;      -- the stop bit is where it should be
                        valid    <= '1';
                    else
                        frame_err <= '1';       -- framing is broken
                        frame     <= (others => '0');
                    end if;
                    state <= STATE_IDLE;        -- either way
                end if;
\end{vhdlcode}

{\itshape (2 marks: 1 for both outcomes and the unconditional return to idle, 1 for the
defaults.)\par}

\medskip
\textbf{What makes the pulses one clock wide.} The defaults are written \textbf{first on every clock
edge}, and only the specific case overrides them. So a \code{valid} set at one edge is cleared at the
next, with no extra counter, no extra state and nothing to get out of step. Writing the default first
and letting the specific case override it is the book's idiom for a single-cycle strobe, and it is
the same one \code{uart_tx} uses for \code{done}.

\textbf{Why the byte is dropped.} \textit{(1 mark)} A framing error means the receiver's idea of
\emph{where the bits were} is wrong: the stop bit was not high where it should have been, so the
eight bits it sampled were sampled at times it cannot vouch for. Delivering them alongside a flag
would put a byte into the RX FIFO that software cannot distinguish from good data once it has popped
past it --- \code{ERROR_FLAGS} is \textbf{global and sticky}, not attached to a FIFO entry, so there
is no way to say ``this byte in particular is suspect''. Dropping it preserves the invariant that
every byte in the RX FIFO is a byte the receiver believes in, and the sticky flag still tells
software that something went wrong.

\textbf{What the bank does with each pulse.} \code{valid} becomes \code{uart_regs}' \code{rx_push},
so the byte enters the RX FIFO and \code{STATUS} bit 1 (RX-valid) goes high, where software can poll
it. \code{frame_err} is latched into \code{ERROR_FLAGS} bit 0, which stays set until software writes
zero to clear it, and \code{STATUS} bit 2 is the OR of the error flags. Both pulses are one clock
wide and neither could be polled directly --- turning them into levels is the register bank's entire
job.

\answerpart{d}{3}
\textbf{All three passed.} \textit{(1 mark)}

\begin{narrowtable}{@{}llll@{}}
  \thead{Byte} & \thead{Binary} & \thead{Reversed} & \thead{Same?} \\
  \midrule
  \code{0xA5} & \code{1010 0101} & \code{1010 0101} & yes \\
  \code{0x3C} & \code{0011 1100} & \code{0011 1100} & yes \\
  \code{0x5A} & \code{0101 1010} & \code{0101 1010} & yes \\
\end{narrowtable}

A receiver that stores its data bits in the reverse order delivers, for each of these,
\textbf{exactly the byte that was sent}. Every check comparing the received byte with the expected
one passed, and no other check in either testbench looks at bit order at all --- so at that revision
the whole of \file{hw/} was blind to it.

\textbf{The property and the count.} They are \textbf{bit-reversal palindromes}. A palindrome is
determined by its top four bits --- \code{b7} fixes \code{b0}, \code{b6} fixes \code{b1}, and so on
--- so there are $2^4 = 16$ of them out of 256. All three constants happen to be among the sixteen.

\textbf{The second, independent reason for \code{uart_top_tb}.} \textit{(1 mark)} It \textbf{loops
\code{tx} back to \code{rx}}. A design built by one person from one set of chapters will get the
order the same way in both halves, and a matched pair of reversals \textbf{cancels exactly}: the
transmitter puts the byte on the wire backwards and the receiver reads it backwards, so the byte read
back over SPI is the byte written, whatever the constant. Changing \code{0x5A} to something
asymmetric would let the testbench catch a \emph{one-sided} flip --- a receiver reversed while the
transmitter is not --- but it structurally cannot catch the matched pair, which is the likelier bug.

\textbf{What the testbench reports today.} \textit{(1 mark)} \code{0x53} is \code{0101 0011} and its
reverse is \code{1100 1010}, so with the reversed-order receiver in place \code{data_out} carries
\textbf{\code{0xCA}}. \code{uart_rx_tb}'s clean-frame case then fails and names both values, since
its message already prints the expected and the received byte through \code{to_hex}:

\begin{console}
uart_rx_tb: received byte mismatch, expected 0x53, got 0xCA!
\end{console}

\noindent The fix that got it there was a constant or two per testbench --- \code{TXBYTE} in
\code{uart_tx_tb}, \code{BYTE_A} and \code{BYTE_B} in \code{uart_rx_tb}, and the transmitted byte
in \code{uart_top_tb} --- which is worth noting in its own right: the gap was a handful of constants
wide and had nothing to do with the checks themselves.

\textbf{What is still out of reach.} A transmitter and a receiver that are reversed
\textbf{together}. \code{uart_tx_tb} and \code{uart_rx_tb} each catch their own half, so the pair
cannot survive both unit testbenches --- but \code{uart_top_tb} alone never could, for the
structural reason above, and a candidate who says ``the loopback is fixed now'' has missed it. No
constant makes a loopback able to test a convention that both ends of it share.

Full marks for reaching the conclusion by any route. The key insight, worth the bulk of the marks, is
that a loopback testbench cannot test a property that both halves get wrong together, and that a
symmetric constant cannot test an ordering.

% ==========================================================================================
\answersection{4}{The bank, written out}{13}

\answerpart{a}{4}
\begin{vhdlcode}
process(clock, reset_s2_n) is
    variable addr: natural range 0 to 15;
begin
    if (reset_s2_n = '0') then
        err_flags <= (others => '0');
    elsif (rising_edge(clock)) then
        addr := to_integer(unsigned(reg_addr));

        -- Software clear: a write of zero to ERROR_FLAGS clears every flag.
        if ((reg_write = '1') and (addr = REG_ERR_FLAGS) and (unsigned(reg_wdata) = 0)) then
            err_flags <= (others => '0');
        end if;

        -- The receiver's pulse sets the framing bit. Written last, so it wins a tie.
        if (frame_err = '1') then
            err_flags(ER_FRAMING) <= '1';
        end if;
    end if;
end process;
\end{vhdlcode}

{\itshape (2 marks: 1 for a reset value of zero with a defined set and a defined clear, 1 for using
\code{uart_def}'s constants and for a clear that is a write of zero to the right index.)\par}

\medskip
\textbf{Both on the same clock.} The set is written \textbf{after} the clear in the same process,
and VHDL gives the last assignment the final say, so \textbf{the set wins}: an error arriving in the
same cycle as a clear is not lost. That is the right way round. Losing an error because software
happened to clear in the same 20\,ns window is a silent data-integrity hole that would be essentially
impossible to reproduce; a flag that survives one extra clear merely makes software clear twice,
which it was polling anyway. A candidate who orders it the other way should say why --- the honest
reason would be ``so a clear always leaves the register readable as zero'', and it is not a good
enough one.

\textbf{Why a pulse cannot be polled.} \textit{(1 mark)} \code{frame_err} is high for \textbf{one
clock}: 20\,ns. Software polls over SPI, and one five-byte transaction at 1\,MHz \code{SCK} takes at
least \textbf{40\,µs} --- two thousand times longer than the pulse, before any inter-byte overhead.
The probability that a \code{STATUS} read lands on precisely the clock the pulse is high is
effectively zero, and even then the pulse would have to coincide with the exact cycle in which the
bridge latches the read value. A pollable bit has to be a \textbf{level that persists until
acknowledged}, and manufacturing that level from a pulse is what the latch is for.

\textbf{Wired straight through.} \textit{(1 mark)} \code{STATUS} bit 2 would read \code{0} on
\textbf{practically every poll, for ever}, no matter how many framing errors occurred: a poll sees it
only if the bridge's one-clock latch happens to fall on the one clock the pulse is high. The system
would appear to have no errors at all.

Worse, it would appear so \emph{consistently}: a driver would never see the bit flicker, never see
it disagree with \code{ERROR_FLAGS}, and have nothing at all to suggest the mechanism was broken. A
flag that is always wrong in the same direction is far harder to notice than one that is sometimes
wrong, and the first evidence would be a bench full of corrupted bytes and a peripheral cheerfully
reporting itself healthy.

\answerpart{b}{4}
\textbf{The read decode.} \textit{(1 mark)}

\begin{vhdlcode}
with to_integer(unsigned(reg_addr)) select reg_rdata <=
    status_reg                           when REG_STATUS,
    ctrl_reg                             when REG_CTRL,
    x"0000" & baud_div_reg               when REG_BAUD_DIV,
    (others => '0')                      when REG_TX_DATA,     -- write-only
    x"000000" & rx_front                 when REG_RX_DATA,
    (others => '0')                      when REG_RX_POP,      -- write-only
    x"0000000" & '0' & err_flags         when REG_ERR_FLAGS,
    (others => '0')                      when others;          -- reserved: 7-15
\end{vhdlcode}

\noindent An equivalent \code{case} inside a combinational process with \code{reg_addr} and every
source register in the sensitivity list is just as good; mark the decode, not the form. What must be
right is that the two write-only registers and every reserved index read \textbf{zero}, and that
every path assigns.

\textbf{The write actions.} \textit{(1 mark)}

\begin{vhdlcode}
tx_push <= '1' when reg_write = '1' and reg_idx = REG_TX_DATA else '0';
rx_pop  <= '1' when reg_write = '1' and reg_idx = REG_RX_POP  else '0';
\end{vhdlcode}

\noindent Two concurrent strobes, outside any process, where \code{reg_idx} is \code{reg_addr}
converted to an integer: \code{tx_push} is the TX FIFO's \code{wr}, with
\code{reg_wdata(7 downto 0)} bound to its \code{wdata}, and \code{rx_pop} is the RX FIFO's
\code{rd}. In the clocked process the two indices do \textbf{nothing}: each FIFO registers its own
strobe on the same edge, and writing FIFO logic into the process as well is the most common way to
get two bytes out of one write. A candidate who registers the strobes inside the process instead,
with defaults of \code{'0'} at the top and \code{'1'} in the two branches, builds a working bank one
clock slower; give the mark if the strobes are provably one clock wide.

Both are \textbf{edge events}: one write, one action, never a held level. That is what makes them
safe to be one-cycle strobes into the FIFOs.

\textbf{Why the read is combinational.} \textit{(1 mark)} The transport contract says the bridge
\textbf{latches the addressed register's value once}, at the end of the command byte, and then
shifts those four bytes out while the master sends dummies. For that latch to capture the right
value, \code{reg_rdata} must already reflect \code{reg_addr} in the same cycle the address becomes
known. A registered read would present its value one clock late, so the bridge would latch the answer
to the \emph{previous} question --- every read would return the previous read's register, which is a
bug that looks like a random off-by-one until somebody notices the pattern.

\textbf{\code{when others}, and why it is a different mistake here.} It must produce
\code{(others => '0')} on the read side and \code{null} on the write side, because the protocol
\textbf{defines} indices 7--15 as reserved: a read returns \code{0x00000000} and a write is ignored.

VHDL does not let you leave the branch out altogether: the choices of a \code{case} or a selected
assignment must cover every value of the selector, and an integer selector has far more than seven,
so a missing \code{others} does not even analyze. What a decoder \emph{can} get wrong is an
\code{others} that assigns nothing --- \code{when others => null;} in a combinational process with no
default assignment above the \code{case} --- and in an ordinary decoder that is a \emph{synthesis}
mistake: it infers a latch, because the output has no value for the uncovered case and must therefore
hold. That is true here too. But here it is also a \textbf{protocol violation}: a driver is entitled
to probe an unimplemented register and get zero, and a bank that returned something else --- or
held the previous value --- would be non-compliant even if it inferred no latch at all. Two
independent reasons for the same line.

\textbf{What the bank relies on.} \textit{(1 mark)} That \code{spi_reg_bridge} asserts
\code{reg_write} \textbf{only on a completed, non-aborted write transaction} (Part~3 of the protocol,
\secref{proto:part3}). A strobe the bank sees is therefore always a real commit. That single
guarantee is what lets \code{TX_DATA} and \code{RX_POP} be simple one-cycle actions with no
commit-and-abort machinery of their own --- the bank never has to think about \code{SS} at all. If a
half-finished transaction could reach it as a \code{reg_write}, every write action in the bank would
need to be held pending and rolled back, and the read/pop split would stop being an elegance and
start being a necessity.

\answerpart{c}{3}
\textbf{The condition.} \textit{(1 mark)}

\begin{vhdlcode}
overrun <= rx_push and rx_full;
\end{vhdlcode}

\noindent The receiver delivered a byte while the RX FIFO had no room, so \code{fifo}'s
\code{not full} guard dropped it. Both signals are already inside \code{uart_regs} ---
\code{rx_push} is the receiver's \code{valid} and \code{rx_full} is the FIFO's own flag --- so the
detection is one AND gate. It sets the flag through the same latch as \code{frame_err}.

\textbf{The bit and the clear path.} \code{ER_OVERRUN}, \textbf{\code{ERROR_FLAGS} bit 2}, cleared
by the same write of \code{0x00000000} to \code{ERROR_FLAGS} that clears the other two, and visible
through \code{STATUS} bit 2 (the OR of the flags) so that a polling driver sees it without an extra
register read.

\textbf{What \code{uart_rx} would need, and why it is worse.} \textit{(1 mark)} It would have to be
told \textbf{whether the previous byte has been consumed} --- in practice the RX FIFO's \code{full}
flag, or a ``byte still unread'' level, as a new input port.

That is a worse design because it inverts the dependency. \code{uart_rx}'s job is to recover bytes
from a wire; it is a purely serial, purely synchronous module with no knowledge of buffering and no
reason to acquire any. Giving it a port that reports the state of a buffer it does not own couples it
to the register bank's internal storage, and means \code{uart_rx_tb} would have to model a FIFO in
order to test a receiver. The bank already holds both facts, one clock apart, in one module.

\textbf{The two physical causes.} \textit{(1 mark)}
\begin{itemize}
  \item \textbf{Overrun}: the wire is fine and the \emph{software} is late. A host busy elsewhere, a
    poll loop that stalled, or simply a line faster than the loop can drain: at 115200 a byte
    arrives every 86.8\,µs, a single five-byte SPI transaction costs at least 40\,µs, and receiving
    one byte takes three of them --- the \code{STATUS} poll, the \code{RX_DATA} read and the
    \code{RX_POP} write --- so at least 120\,µs. A continuous stream at 115200 outruns even a driver
    that does nothing else, and the eight-entry FIFO only postpones the loss. It is a timing failure
    above the peripheral.
  \item \textbf{Framing error}: the \emph{wire} is not fine. A baud mismatch large enough to move the
    stop-bit sample out of its bit, a break from a transmitter that lost power or a cable that came
    loose, noise or a reflection on a long unterminated run, or a level-shifting problem. It is an
    electrical or configuration failure below the peripheral.
\end{itemize}
One is caused by the reader, the other by the sender or the medium, which is exactly why they are
separate bits rather than one ``something went wrong'' flag.

\answerpart{d}{2}
\textbf{Two disagreements with \code{uart_regs_tb}.} \textit{(1 mark each, any two of the
following)}
\begin{enumerate}
  \item \textbf{After reset.} The testbench checks that \code{STATUS} immediately after reset shows
    \textbf{TX-ready and TX-idle set} with RX-valid and Error clear. A register updated only when
    \code{rx_push}, \code{tx_pop} or \code{frame_err} asserts has seen none of those, so unless its
    reset value is hand-written to \code{0b1001} it reads zero and the check fails --- and a
    hand-written reset value is a second copy of the truth to keep in step with the FIFOs for ever.
  \item \textbf{\code{RX_POP}.} Popping the RX FIFO empties it and must clear RX-valid, and the
    testbench checks exactly that. But \code{RX_POP} is a \textbf{register write}, not one of the
    three events listed, so RX-valid stays set after the FIFO is empty. A driver would then poll, see
    RX-valid, read \code{RX_DATA} from an empty FIFO and pop it, for ever.
  \item \textbf{\code{TX_DATA}.} A write pushes a byte and must clear TX-idle; a push is also not in
    the list, so the testbench's ``a \code{TX_DATA} write must clear TX-idle'' check fails for the
    same reason.
  \item \textbf{\code{tx_busy}.} It changes with no event at all --- the transmitter simply enters
    and leaves a state --- so TX-idle goes stale whenever the datapath moves without the bank being
    told.
\end{enumerate}

\textbf{The rule.} Derive a status bit from \textbf{the state it reports}, never from the events that
changed that state. A register written by events is wrong whenever an event is missed, whenever two
coincide, whenever something changes without an event, and immediately after reset.

% ==========================================================================================
\answersection{5}{Writing the contracts}{12}

\answerpart{a}{5}
\begin{cppcode}
/**
 * @file Byte transport interface.
 */
#pragma once

#include <stdint.h>

namespace driver
{
namespace transport
{
/**
 * @brief Byte-level SPI transport interface.
 */
class Interface
{
public:
    /** @brief Destructor. */
    virtual ~Interface() noexcept = default;

    /** @brief Start a transaction by pulling SS low. */
    virtual void begin() noexcept = 0;

    /**
     * @brief Exchange a single byte in full duplex.
     *
     * @param[in] byte Byte to send.
     * @return The byte received at the same time.
     */
    virtual uint8_t transfer(uint8_t byte) noexcept = 0;

    /** @brief End the transaction by releasing SS. */
    virtual void end() noexcept = 0;
};
} // namespace transport
} // namespace driver
\end{cppcode}

{\itshape (2 marks: 1 for the three pure virtuals with correct signatures, 1 for the
\textbf{virtual}, \code{noexcept}, \code{= default} destructor and the nested-block namespaces.)\par}

\begin{cppcode}
/**
 * @file UART driver interface.
 */
#pragma once

#include <stdint.h>

namespace driver
{
namespace uart
{
/**
 * @brief UART driver interface.
 */
class Interface
{
public:
    /** @brief Destructor. */
    virtual ~Interface() noexcept = default;

    /** @brief Configure the peripheral: set the baud divider and enable it. */
    virtual void configure(uint16_t baudDiv) noexcept = 0;

    /** @brief Send one byte. @return True if accepted, false if the TX FIFO was full. */
    virtual bool write(uint8_t byte) noexcept = 0;

    /** @brief Receive one byte. @return True if a byte was retrieved, false otherwise. */
    virtual bool read(uint8_t& byte) noexcept = 0;

    /** @brief Read the raw STATUS register. */
    virtual uint32_t status() const noexcept = 0;

    /** @brief Read the raw ERROR_FLAGS register. */
    virtual uint32_t errorFlags() const noexcept = 0;

    /** @brief Clear the error flags. */
    virtual void clearErrors() noexcept = 0;
};
} // namespace uart
} // namespace driver
\end{cppcode}

{\itshape (3 marks: 1 for \code{configure}, \code{write} and \code{read} with the right parameter and
return types --- note that \code{read} takes a \textbf{reference} out-parameter and returns
\code{bool}; 1 for the three status methods with \code{status()} and \code{errorFlags()}
\textbf{\code{const}} and \code{clearErrors()} not; 1 for \code{noexcept} throughout and for the
AVR-portable style: \code{<stdint.h>}, bare types, nested blocks, no \code{[[nodiscard]]}.)\par}

\medskip
Deduct for \code{std::uint8_t}, for \code{namespace driver::uart}, for \code{[[nodiscard]]}, for a
non-virtual destructor, and for \code{read} returning the byte instead of a \code{bool}.

\answerpart{b}{3}
\begin{cppcode}
namespace reg
{
/** UART status register. */
constexpr uint8_t STATUS{0U};

/** UART control register. */
constexpr uint8_t CTRL{1U};

/** UART baud rate divider register. */
constexpr uint8_t BAUD_DIV{2U};

/** UART TX data register. */
constexpr uint8_t TX_DATA{3U};

/** UART RX data register. */
constexpr uint8_t RX_DATA{4U};

/** UART RX pop register. */
constexpr uint8_t RX_POP{5U};

/** UART error flag register. */
constexpr uint8_t ERROR_FLAGS{6U};
} // namespace reg

namespace status
{
/** Transmitter is ready for another byte. */
constexpr uint8_t TX_READY{0U};

/** An RX byte is available. */
constexpr uint8_t RX_VALID{1U};

/** UART is in an error state. */
constexpr uint8_t ERROR{2U};

/** Transmitter is idle. */
constexpr uint8_t TX_IDLE{3U};
} // namespace status
\end{cppcode}

{\itshape (1 mark for the seven indices, 1 for the four status positions --- both sets in the right
order with the right values.)\par}

\medskip
\textbf{Why plain \code{constexpr}.} \textit{(1 mark, shared with the next point)} \code{inline}
\textbf{variables} are a C++17 feature, and the oldest AVR toolchain the course targets, the avr-gcc
Microchip Studio ships, does not have them --- it does not accept \code{-std=c++17} at all. Plain
\code{constexpr} compiles everywhere the course needs it to.

\textbf{What internal linkage costs here: nothing.} At namespace scope \code{constexpr} already
implies internal linkage, so each translation unit gets its own copy of a \textbf{compile-time}
constant. That copy is folded into the instruction stream at every use site and occupies no storage,
so there is no duplication to pay for and no one-definition-rule problem, because there is no object
to define twice.

\textbf{Positions rather than masks.} It keeps the C++ \textbf{identical to \file{uart_def.vhd}},
which stores \code{natural} bit positions because that is what indexes a \code{std_logic_vector},
and identical to the protocol's own ``bit N'' wording. The two sides are then each checked against
the specification rather than against each other, which is the only arrangement in which a
disagreement is findable. The mask is formed at the use site:

\begin{cppcode}
if (status & (1U << status::RX_VALID)) { /* a byte is waiting */ }
\end{cppcode}

\noindent which is explicit about which bit is meant and cannot be mistaken for a value.

\answerpart{c}{2}
\textbf{What it is for, and who owns it.} \textit{(1 mark)} \code{myStop} refers to a flag the
\textbf{caller} owns --- the very flag the application's \code{run(const volatile bool& stop)} is watching.
The stub sets it \code{true} the moment its scripted RX buffer is exhausted, so a single-threaded
test can end an application loop that would otherwise never return.

\code{read()} therefore does two things when the buffer runs out: it sets \code{myStop} to
\code{true} and returns \code{false}. Both matter --- the \code{false} says ``no byte this pass'',
the flag says ``and there will not be another''.

\textbf{Why a stub should care at all.} Because the thing it doubles for is a \textbf{UART}, and a
real UART never runs out of input; it merely has none \emph{yet}. There is no in-band way to say
``that was all the input there will ever be'', so the double supplies an out-of-band one. It is the
double being faithful about the test's needs rather than about the hardware, which is exactly what a
double is for.

\textbf{What the reference member forces.} \textit{(1 mark)} The language forces two things. The
\textbf{default constructor} is implicitly deleted, because a reference member must be bound at
construction --- there is no way to build a \code{Stub} without a flag to point at. And the copy and
move \textbf{assignment} operators are implicitly deleted, because a reference cannot be rebound
after construction. The copy and move \textbf{constructors} are \emph{not} forced: the implicit ones
would compile, and bind the new stub to the same flag. The book deletes them anyway, explicitly,
because a copy that shares its original's stop flag but not its buffers is a trap, and a class that
cannot be assigned but can be copied invites exactly that mistake. \code{Uart} and \code{EchoNode}
delete the same five for the same reason. Award the mark for saying which one is forced, and why the
other two are deleted by choice.

\answerpart{d}{2}
\textbf{The rule.} \textit{(1 mark)} A test double must be faithful about \textbf{everything the
code under test depends on}, and may ignore everything else. It is a fixture, not a model of the
hardware: fidelity it is not asked for is fidelity nobody checks, and unchecked behaviour in a double
is a liability the day it disagrees with the real thing.

\textbf{Applied here.} The only application tested against \code{driver::uart::Stub} is
\code{app::EchoNode}, which drives the UART entirely through \code{read()} and \code{write()}. It
never calls \code{status()}, \code{errorFlags()} or \code{clearErrors()}. So returning \code{0} and
doing nothing is not laziness --- it is the correct amount of fidelity, and inventing a status word
would mean shipping behaviour with nothing to check it against.

\textbf{What would have to change.} \textit{(1 mark)} The moment an application makes a
\textbf{decision} from the status word, the stub needs a scriptable one: a settable \code{uint32_t},
a helper such as \code{setStatus(uint32_t)} and \code{setErrorFlags(uint32_t)}, and a
\code{clearErrors()} that actually zeroes what \code{errorFlags()} returns so that the
clear-then-recheck path can be exercised.

Examples of such a thing: a node that stops or reports when \code{STATUS} bit 2 goes high; one that
waits for TX-idle before letting the MCU sleep; one that counts framing errors and resets the link
after a threshold. Each of those is a decision the current stub cannot script, and each is a good
reason to extend it --- but only when it exists.

% ==========================================================================================
\answersection{6}{Writing the protocol}{13}

\answerpart{a}{5}
\begin{cppcode}
// -----------------------------------------------------------------------------
uint32_t Uart::readReg(const uint8_t addr) const noexcept
{
    constexpr uint8_t dummy{0x00U};
    constexpr uint32_t shift{8U};
    constexpr uint32_t size{static_cast<uint32_t>(sizeof(uint32_t))};
    constexpr uint32_t last{size - 1U};

    uint32_t value{};
    auto& transport = const_cast<transport::Interface&>(myTransport);

    transport.begin();
    transport.transfer(addr);              // bit 7 clear => read. Reply ignored.

    for (uint8_t i{}; i < size; ++i)
    {
        const auto idx = last - i;         // 3, 2, 1, 0 => most significant byte first
        const uint8_t response{transport.transfer(dummy)};
        value |= static_cast<uint32_t>(response) << (shift * idx);
    }
    transport.end();
    return value;
}

// -----------------------------------------------------------------------------
void Uart::writeReg(const uint8_t addr, const uint32_t value) noexcept
{
    constexpr uint8_t writeMask{0x80U};
    constexpr uint32_t shift{8U};
    constexpr uint32_t size{static_cast<uint32_t>(sizeof(uint32_t))};
    constexpr uint32_t last{size - 1U};

    myTransport.begin();
    myTransport.transfer(static_cast<uint8_t>(addr | writeMask));   // bit 7 set => write

    for (uint32_t i{}; i < size; ++i)
    {
        const auto idx = last - i;         // most significant byte first
        myTransport.transfer(static_cast<uint8_t>(value >> (shift * idx)));
    }
    myTransport.end();
}
\end{cppcode}

{\itshape (4 marks: 1 for the command bytes --- the raw index for a read, \code{index | 0x80} for a
write; 1 for the \code{begin()}, five transfers, \code{end()} framing in both; 1 for \textbf{most
significant byte first} in the read assembly; 1 for \textbf{most significant byte first} in the
write split.)\par}

\medskip
Byte order is the part to be strict about. Reverse either loop and the same four bytes cross the
wire carrying a different 32-bit value, and neither the compiler nor a round trip through a
symmetrically-wrong stub will say a word.

\textbf{The replies.} \code{readReg} \textbf{discards} the reply to the command byte: the slave was
shifting that out before it knew which register was being asked for, so it is meaningless by
construction. \code{writeReg} discards \textbf{all four} of its replies as well --- on a write the
master has nothing to learn from \code{MISO}, and the transfers happen only because SPI is duplex and
a byte out costs a byte in. \textit{(Half a mark.)}

\textbf{What \code{const} buys.} \textit{(Half a mark.)} It lets \code{status()} and
\code{errorFlags()} --- the two public methods that are pure observations --- be \code{const} too.
Those are the two that depend on it. Without it, a caller holding a \code{const Uart&} could not
read the peripheral's status at all, which would be an odd thing for a \code{const} handle to forbid.

\answerpart{b}{4}
\begin{cppcode}
// -----------------------------------------------------------------------------
bool Uart::read(uint8_t& byte) noexcept
{
    const auto status  = readReg(reg::STATUS);
    const auto rxValid = static_cast<bool>(status & (1U << status::RX_VALID));

    if (rxValid)
    {
        constexpr uint32_t popCount{1U};
        byte = readReg(reg::RX_DATA);        // pure read: does not pop
        writeReg(reg::RX_POP, popCount);     // separate write: advances the FIFO
    }
    return rxValid;
}
\end{cppcode}

{\itshape (1 mark: poll, then read, then pop, in that order, with nothing issued when
\code{RX_VALID} is clear.)\par}

\medskip
\textbf{The recorded bytes.} \textit{(1 mark)}

\begin{console}
transaction 1  STATUS read    : 00 00 00 00 00
transaction 2  RX_DATA read   : 04 00 00 00 00
transaction 3  RX_POP write   : 85 00 00 00 01
\end{console}

\noindent\code{beginCalls() == 3}, \code{endCalls() == 3}, and \code{txLen() == 15}.

\textbf{The scripting.} \textit{(1 mark)}

\begin{cppcode}
stub.injectRxByte(0x00U);                    // command-phase placeholder
stub.injectRxWord(1U << status::RX_VALID);   // STATUS = 0x00000002
stub.injectRxByte(0x00U);                    // command-phase placeholder
stub.injectRxWord(0x00000041U);              // RX_DATA = 0x41
\end{cppcode}

\noindent\textbf{Ten bytes} in total, for \textbf{two} register values. \textit{(1 mark)}

\textbf{The discrepancy.} The stub is a \textbf{byte pipe}: it returns the next queued byte on
\emph{every} \code{transfer()} call, and a five-byte read transaction makes five of them. The reply
the driver actually uses is only four bytes long, because the reply to the command byte is discarded
--- but it is still \emph{consumed}. So each read costs one placeholder plus the four data bytes.

The third transaction needs nothing queued: it is a write, the driver ignores all five replies, and
the stub returns \code{0x00} once its buffer is exhausted anyway.

Getting this wrong is spectacular and confusing: queue only the four data bytes and every read is
shifted by one byte. The command phase swallows the most significant byte, so \code{readReg}
assembles the value's low three bytes, shifted up by eight bits, together with the first byte queued
for the \emph{next} read: \code{STATUS} reads back as \code{0x00000200}, RX-valid is clear, and
\code{read()} returns \code{false} with a byte waiting. Every later read in the test is misaligned
the same way. This is precisely why the provided suite wraps the scripting in a
\code{scriptRead(stub, value)} helper rather than leaving it at each call site.

\answerpart{c}{2}
\begin{cppcode}
inline void writeBlocking(Interface& uart, const uint8_t byte) noexcept
{
    while (!uart.write(byte)) {}
}

inline void readBlocking(Interface& uart, uint8_t& data) noexcept
{
    while (!uart.read(data)) {}
}
\end{cppcode}

{\itshape (1 mark for both, \code{inline} and \code{noexcept}, in the \code{driver::uart}
namespace.)\par}

\medskip
\textbf{Why \code{inline}.} \textit{(1 mark, shared)} They are \textbf{defined in a header}, so every
translation unit that includes it emits a definition. Without \code{inline} that is a
one-definition-rule violation and a duplicate-symbol error at link time the moment two files include
the header.

\textbf{Why \code{Interface&} and not \code{Uart&}.} So they work over \textbf{any} implementation
--- the concrete \code{Uart}, the \code{driver::uart::Stub} of \chapref{6}, any future driver ---
dispatching virtually at run time. Taking the concrete type would tie a convenience function to one
implementation for no reason.

\textbf{Why free functions in a separate header.} They add \textbf{no state and no new contract}:
they are a spin loop over an operation the interface already provides. Putting them in the interface
would oblige every implementation, including every stub, to reimplement that loop, and would push the
decision to block \emph{below} the seam, so a test could no longer opt out of it. The core stays
deterministic and testable, and the spinning lives where a caller chooses it.

\answerpart{d}{2}
\textbf{The call sequence.} \textit{(1 mark)}

\begin{console}
status()  ->  readReg(STATUS)  ->  write(cmd)  ->  readReg(STATUS)  ->  write(cmd)  ->  ...
\end{console}

\noindent\code{Uart::write} polls \code{STATUS} before it sends anything, and \code{readReg} is what
polls. So a \code{readReg} implemented in terms of \code{write()} calls the very method that calls
it, with \textbf{no base case}: infinite mutual recursion, entered the first time anything reads a
register: the first \code{read()}, \code{write()} or \code{status()} after \code{configure()}, which
only writes, and every one thereafter.

\textbf{What the ATmega328P exhibits.} \textit{(1 mark)} Each level pushes a stack frame. The
ATmega328P has \textbf{2\,KB of SRAM}, the stack grows down from the top of it toward the statically
allocated data, and there is no MMU, no stack guard, no exception and no fault handler. So the stack
simply runs into \code{.data} and \code{.bss} and starts overwriting live variables, then keeps
going.

The symptom is not a clean crash: it is \textbf{memory corruption followed by a jump to a garbage
address}, which typically presents as a device that appears to reset in a loop the moment the driver
is first used, or that behaves erratically before it does. On the host the same bug is a
segmentation fault, which is the kinder outcome and the reason the host suite finds it first --- one
more argument for building the driver on the host before it ever reaches a chip.

\textbf{The rule.} \code{transfer()} moves a \textbf{byte on the wire}; \code{write()} sends a
\textbf{byte over the UART}. They are two different layers that happen to have the same shape, and
the private register core may only ever touch \code{myTransport}.

% ==========================================================================================
\answersection{7}{Writing the transport}{11}

\answerpart{a}{5}
\begin{cppcode}
// -----------------------------------------------------------------------------
AvrSpi::AvrSpi() noexcept
{
    DDRB |= (1U << SCK) | (1U << MOSI) | (1U << SS);   // SCK, MOSI, SS out; MISO stays in.
    PORTB |= (1U << SS);                               // Idle the chip select high.
    SPCR = (1U << SPE) | (1U << MSTR) | (1U << SPR0);  // Enable, master, mode 0, MSB, f_osc/16.
}

// -----------------------------------------------------------------------------
AvrSpi::~AvrSpi() noexcept
{
    DDRB &= ~((1U << SCK) | (1U << MOSI) | (1U << SS));
    PORTB &= ~(1U << SS);
    SPCR = 0U;
}

// -----------------------------------------------------------------------------
void AvrSpi::begin() noexcept { PORTB &= ~(1U << SS); }

// -----------------------------------------------------------------------------
uint8_t AvrSpi::transfer(const uint8_t byte) noexcept
{
    SPDR = byte;
    while (0U == (SPSR & (1U << SPIF))) {}   // Spin until the transfer completes.
    return SPDR;                             // Now SPDR holds the received byte.
}

// -----------------------------------------------------------------------------
void AvrSpi::end() noexcept { PORTB |= (1U << SS); }
\end{cppcode}

{\itshape (3 marks: 1 for the constructor's three statements with \code{SS} idled high and
\code{MISO} untouched; 1 for \code{begin()} low and \code{end()} high; 1 for \code{transfer()} with
the \code{SPIF} spin between the write and the read.)\par}

\medskip
\code{SPR0} set with \code{SPR1} clear and \code{SPI2X} clear gives f\textsubscript{osc}/16, which is
1\,MHz. \code{DORD}, \code{CPOL} and \code{CPHA} are all clear, which is MSB first, mode 0 --- so the
contract is met by bits that are \emph{absent} as much as by bits that are present.

\textbf{The destructor.} \textit{(1 mark)} It must clear \textbf{exactly what the constructor set
and nothing else}:
\begin{itemize}
  \item the three \code{DDRB} direction bits for \code{SCK}, \code{MOSI} and \code{SS}, putting
    those pins back to inputs;
  \item the \code{SS} bit in \code{PORTB}, which was the chip-select drive;
  \item all of \code{SPCR}, written to \code{0}, which disables the peripheral.
\end{itemize}
It must \textbf{leave alone}: \code{MISO}, because the constructor never touched it --- it is an
input at reset and stays one, and returning something you did not take is not RAII, it is trespass.
Also \code{SPSR}, which the constructor never wrote: its one configuration bit, \code{SPI2X}, is left
at its reset value of 0, which f\textsubscript{osc}/16 relies on. And every other bit of \code{DDRB}
and \code{PORTB}, which belong to whatever else is using port B. The \code{&=} and \code{|=} forms
are what confine the edits to those bits.

\textbf{Why \code{SPCR = ...} and not \code{SPCR |= ...}.} \textit{(1 mark)} The plain assignment
\textbf{guarantees} that every bit the design relies on being clear is clear: \code{DORD},
\code{CPOL}, \code{CPHA}, \code{SPIE} and \code{SPR1}. With \code{|=}, whatever was in \code{SPCR}
beforehand survives --- anything that ran first, an Arduino core's \code{SPI.begin()} or a library
of your own, could have left \code{DORD} set (LSB first) or \code{SPR1} set (a different prescaler),
and the transport would silently run at the wrong bit order or the wrong rate with no line in the
source to blame it on.

The general rule, worth stating: a register you \textbf{own} is \emph{set}; a register you
\textbf{share} is \emph{merged}. \code{SPCR} is entirely this class's, so it is assigned.
\code{DDRB} and \code{PORTB} are shared with every other pin on port B, so they are
read-modify-written.

\answerpart{b}{3}
\textbf{The principle.} \textit{(1 mark)} \textbf{RAII}: what the constructor acquires, the
destructor releases. This class acquires the SPI hardware and three port pins, and a defaulted
destructor gives none of them back. The \code{Interface} declares a \textbf{virtual} destructor
precisely so that destroying an \code{AvrSpi} through an \code{Interface&} runs the derived one;
providing that mechanism and then making it do nothing throws away the thing the base class went to
the trouble of arranging.

\textbf{What is never given back.} PB5, PB3 and PB2 are left as \textbf{outputs} in \code{DDRB};
PB2 is left \textbf{driven high} in \code{PORTB}; and \code{SPCR} still has \code{SPE} and
\code{MSTR} set, so the peripheral is still enabled and still owns those three pins.

\textbf{Concrete harm.} \textit{(1 mark)} Any code that runs afterwards and wants PB2, PB3 or PB5 as
a general-purpose \textbf{input} finds them driven as outputs and reads its own drive rather than the
world. Any code that wants the SPI peripheral in a different configuration --- as a slave, at a
different prescaler, LSB first --- inherits a half-configured master rather than the reset state, and
a transport that uses \code{|=} on \code{SPCR} (see (a)) would inherit \code{MSTR} and never notice.
And with \code{SPE} still set the peripheral keeps driving \code{SCK} and \code{MOSI}, so anything
else trying to use those pins fights the hardware for them.

\textbf{Why ``never destroyed'' is unsafe here in particular.} \textit{(1 mark)} Because this
codebase builds its objects as \textbf{automatic variables and injects them by reference}, and the
\textbf{host test suite} does exactly that: it constructs an \code{AvrSpi} inside a scope, exercises
it over the mocked register file, lets it go out of scope, and constructs another for the next case.
With a do-nothing destructor the peripheral it configured outlives it, so a defaulted destructor is
observably wrong even on the host. The suite defends itself today by calling \code{resetHardware()}
at the top of every case; without that defence the cases would become \textbf{order-dependent} and
the configuration assertions could pass for the wrong reason --- the bits right because the last
test set them, not because this constructor did. The point is that a class should not need its tests
to clean up after it.

\answerpart{c}{3}
\textbf{What sets each.} \textit{(1 mark)}
\begin{itemize}
  \item \textbf{\code{SCK}} is set by the SPI \textbf{prescaler}, a hardware divider off the CPU
    clock selected by \code{SPR1:SPR0} and \code{SPI2X}: f\textsubscript{osc}/4, /16, /64, /128 and
    the doubled variants. It divides whatever the chip is \emph{actually} running at, and no software
    constant takes part --- which is exactly why \code{AvrSpi} needs no \code{F_CPU}.
  \item \textbf{The USART's baud} is set by \textbf{\code{UBRR0}, a number the software computes}:
    \code{UBRR0 = F_CPU / (16 * baud) - 1}. The compiler substitutes \code{F_CPU} at compile time,
    so the number is only right if the macro matches reality. The logging path needs \code{F_CPU}
    because it does arithmetic with it; the transport does not because it does none.
\end{itemize}

\textbf{At 8\,MHz with \code{F_CPU} still 16000000.} \textit{(2 marks)}
\begin{itemize}
  \item \textbf{\code{SCK}} becomes 8\,MHz\,/\,16 = \textbf{500\,kHz} instead of 1\,MHz. The link
    \textbf{still works}: the FPGA synchronizes and edge-detects \code{sclk} into its own 50\,MHz
    domain, so it keeps up with far more than 1\,MHz and there is no minimum rate. Everything is
    simply twice as slow --- a five-byte transaction goes from 40\,µs to 80\,µs.
  \item \textbf{The USART} computes \code{UBRR0} for a 16\,MHz clock while the hardware divides a
    real 8\,MHz one, so the actual baud is \textbf{half} of what the divider was computed for: at a
    115200 log rate, about 57600. A terminal set to 115200 sees framing errors and garbage, with the
    occasional plausible character by luck.
\end{itemize}

\textbf{Which you notice first, and why that is unfortunate.} The \textbf{terminal}, immediately and
unmistakably --- garbage on a console is impossible to miss, while a \code{readReg} that takes
80\,µs instead of 40\,µs looks exactly like one that takes 40\,µs unless you are measuring.

That is unfortunate because the loud symptom points at the \textbf{logging}, the one part of the
system that does not matter, and hides the fact that the SPI link is also running at half speed ---
which is precisely what Exercise~\ref{c9:ex:3} asks you to measure against the 1\,MHz budget. A
candidate who ``fixes'' it by halving the log baud has made the symptom go away and left the wrong
clock in place, and every timing figure they take afterwards is out by a factor of two. The fix is to
\code{F_CPU}, so that it states the clock the part really runs at (or to the fuses, if 16\,MHz was
what was wanted).

% ==========================================================================================
\answersection{8}{The application, and the plan}{12}

\answerpart{a}{4}
\begin{cppcode}
/**
 * @file Echo node implementation.
 */
#pragma once

#include "app/interface.hpp"

// clang-format off
namespace driver
{
/** UART driver interface. */
namespace uart { class Interface; }
} // namespace driver
// clang-format on

namespace app
{
/**
 * @brief Echo node implementation.
 *
 *        This class is non-copyable and non-movable.
 */
class EchoNode final : public Interface
{
public:
    /** @brief Constructor. @param[in] uart UART driver instance. */
    explicit EchoNode(driver::uart::Interface& uart) noexcept;

    /** @brief Destructor. */
    ~EchoNode() noexcept override = default;

    /** @brief Run application. @param[in] stop Set true to stop the application. */
    void run(const volatile bool& stop) noexcept override;

    EchoNode()                           = delete;
    EchoNode(const EchoNode&)            = delete;
    EchoNode(EchoNode&&)                 = delete;
    EchoNode& operator=(const EchoNode&) = delete;
    EchoNode& operator=(EchoNode&&)      = delete;

private:
    /** UART driver instance. */
    driver::uart::Interface& myUart;
};
} // namespace app
\end{cppcode}

\begin{cppcode}
/**
 * @file Echo node implementation details.
 */
#include "app/echo_node.hpp"
#include "driver/uart/blocking.hpp"
#include "driver/uart/interface.hpp"

namespace app
{
// -----------------------------------------------------------------------------
EchoNode::EchoNode(driver::uart::Interface& uart) noexcept
    : myUart{uart}
{}

// -----------------------------------------------------------------------------
void EchoNode::run(const volatile bool& stop) noexcept
{
    while (!stop)
    {
        uint8_t rxByte{};

        if (myUart.read(rxByte)) { driver::uart::writeBlocking(myUart, rxByte); }
    }
}
} // namespace app
\end{cppcode}

{\itshape (3 marks: 1 for the reference member to the \textbf{interface} with the five deleted
operations; 1 for the \code{explicit}, \code{noexcept} constructor that does no I/O and the
\code{override} on \code{run}; 1 for a loop that re-checks \code{stop} every pass and echoes only on
a successful \code{read}.)\par}

\medskip
\textbf{Why non-blocking read, blocking write.} \textit{(1 mark)} A blocking read would spin inside
\code{Interface::read()} until a byte arrived and would never return to the top of the loop, so
\code{stop} would never be read again --- the application could not be stopped by anybody, including
its own test, which would hang. Polling \code{read()} returns every pass, which is the only thing
that gives the flag a chance to be seen.

The echo write may block because the byte is \textbf{already in hand}: there is no decision left to
make, nothing else the loop could usefully do, and no deadlock available, since the TX FIFO is
drained by hardware whether software cooperates or not. Poll where you might have to give up; block
where you have already committed.

\textbf{Why a \code{volatile bool} by \code{const} reference.} Freestanding avr-libc ships no
\code{<atomic>}, so \code{std::atomic<bool>} is not available --- and on the single-core ATmega328P a
byte read is already indivisible, so it would not buy anything. It is \code{volatile} so that every
pass reads it from memory: \code{run()} never writes it, and when it is set from outside anything the
loop calls, by an interrupt handler for instance, a plain \code{bool} could be read once and kept in
a register. It is passed \textbf{by reference}
and re-read every pass so that the owner can set it at any moment; a return value would only be
examined when \code{run()} returned, which is exactly what it will not do until the flag is set. And
it is \code{const} because the application only ever \textbf{reads} it: setting it is the caller's
business.

\answerpart{b}{3}
\textbf{The test.} \textit{(2 marks)}

It runs over the \textbf{\code{driver::uart::Stub}} of \chapref{6} --- the UART-level double, not the
transport stub, since \code{EchoNode} sits above the driver and knows nothing of SPI. Host test code
may use full modern C++.

\begin{cppcode}
volatile bool stop{false};
driver::uart::Stub stub{stop};
app::EchoNode node{stub};

stub.injectRxByte(0x00U);
stub.injectRxByte(0x41U);
stub.injectRxByte(0xFFU);

node.run(stop);

EXPECT_EQ(stub.txLen(), 3U);
EXPECT_EQ(stub.txBuf()[0], 0x00U);
EXPECT_EQ(stub.txBuf()[1], 0x41U);
EXPECT_EQ(stub.txBuf()[2], 0xFFU);
\end{cppcode}

\begin{itemize}
  \item \textbf{What it queues:} the scripted RX bytes, \code{0x00}, \code{0x41} and \code{0xFF}. The
    two extremes are chosen deliberately --- they are the values most likely to expose a truncation,
    a sign-extension or a ``treat zero as no data'' mistake.
  \item \textbf{How the loop terminates:} the stub sets \code{stop} to \code{true} when \code{read()}
    finds its RX buffer exhausted, and \code{run()} sees it on the next pass. That is also a
    \textbf{check on the implementation}: because termination depends on \code{run()} polling and
    re-checking, a \code{run()} that blocked in \code{readBlocking()} would hang the test rather than
    fail it --- the test telling you the loop is not actually stoppable.
  \item \textbf{What it asserts:} the recorded TX length equals the number queued, and each recorded
    byte equals the queued byte at the same position --- equality \textbf{in order}, because a UART
    is a byte stream and order is half the contract.
\end{itemize}

\textbf{The empty case.} \textit{(1 mark)} Queue nothing, call \code{run(stop)}, and assert that it
returns and that \code{txLen()} is \code{0}.

It proves two things: that the loop's exit does \textbf{not} depend on having echoed anything, and
that the first thing the loop does before any I/O is test \code{stop} --- so a caller who sets the
flag before \code{run()} is even entered gets an immediate return.

The implementation mistake it isolates is a \code{run()} that \textbf{ignores \code{read()}'s return
value} and echoes \code{rxByte} regardless. With nothing queued, that is the only thing the loop can
do wrong, and it shows as one spurious \code{0x00} where there should be none. The three-byte case
catches it too, but only through its length assertion: every byte it compares is correct, and the bug
adds a fourth, spurious \code{0x00} after them, so \code{txLen()} reads 4 instead of 3. A test that
checked only the first \emph{n} bytes would pass it, which is why the length is asserted as well as
the contents.

\answerpart{c}{3}
\textbf{The five rungs, and what each adds.} \textit{(1 mark)}

\begin{booktable}{@{}p{2.2em}p{11.2em}L@{}}
  \thead{Rung} & \thead{What it is} & \thead{The one layer it adds} \\
  \midrule
  a & Data-plane \textbf{pin loopback}: the wrapper wires \code{rx} straight back to \code{tx},
    bypassing \code{uart_top} & The USB-serial adapter, its 3.3\,V levels, the terminal's baud and
    frame settings, the two data-plane wires and their crossover, and the FPGA pin assignment.
    \textbf{None of your logic.} \\
  b & \textbf{The control plane}: write \code{BAUD_DIV} over SPI and read it back & The four SPI
    lines and the level shifter, \code{AvrSpi}, the provided \code{spi_slave} and
    \code{spi_reg_bridge}, and the register bank's write and read paths. \\
  c & \textbf{Peripheral loopback}: \code{tx} tied to \code{rx} on the board & The whole datapath
    on real silicon at a real 50\,MHz --- \code{baud_gen}, \code{uart_tx}, \code{uart_rx}, both
    FIFOs and the TX feeder. \\
  d & \textbf{The real data plane}: against the terminal & A \textbf{second device} with an
    \textbf{independently generated baud rate} that must now agree with the peripheral's own. \\
  e & \textbf{\code{app::EchoNode}} & The application. \\
\end{booktable}

\textbf{Why the order is forced.} \textit{(1 mark)} \textbf{\code{BAUD_DIV}.} The peripheral cannot
transmit or receive anything usable until a divider has been written to it --- \code{baud_gen} has
no default and the guarded conversion substitutes \code{1}, which is 3.125\,Mbaud and useless. And
the \textbf{only} route to \code{BAUD_DIV} is an SPI transaction: it is not reachable from the data
plane, not from a pin, not from a strap, not from its reset value, which is zero.

So the control plane must be proven before \textbf{any} rung that uses the peripheral. Any ladder
that puts a peripheral rung first is exercising two layers at once and calling it one, and its first
failure names both.

\textbf{Two reasons EchoNode-first is worse even when it succeeds.} \textit{(1 mark)}
\begin{enumerate}
  \item \textbf{A pass produces no record.} It tells you the whole stack works right now; it does not
    tell you which layers you may trust tomorrow. The next change --- a re-synthesis, a different
    adapter, a wire moved, a new bench --- puts you back to a single binary result with no cleared
    layers to fall back on. The ladder's product is not the final pass, it is the sequence of things
    each rung ruled out, and one end-to-end test produces none of them.
  \item \textbf{It hides marginal failures.} The rungs are designed to stress \emph{different}
    things: rung a exercises the levels and the terminal settings with none of your logic in the
    path, rung c removes the second baud source entirely, rung d puts it back. A stack that echoes
    correctly today may still be running with a 0.3\,V \code{MISO} margin, an unshifted line that
    happens to work at this temperature, or a baud error that is inside tolerance now and outside it
    later. Each rung would have failed on exactly one of those; a passing end-to-end test reports
    none.
\end{enumerate}
(The usual argument --- that it is also worse when it \emph{fails}, because nothing is localized ---
is correct but was not what the question asked.)

\answerpart{d}{2}
\textbf{What is now impossible.} \textit{(1 mark)} \textbf{Host-testing \code{EchoNode} without a
transport.} A \code{Uart} can only be constructed over a \code{driver::transport::Interface&}, so
the test must build a \code{driver::transport::Stub}, then a \code{Uart} over it, then script
\textbf{every SPI transaction the echo produces}: for each byte received, a \code{STATUS} read, an
\code{RX_DATA} read and an \code{RX_POP} write; for each byte echoed, a \code{STATUS} read and a
\code{TX_DATA} write --- five bytes each, and each read needing its command-phase placeholder plus
four data bytes queued. Every transfer consumes a queued byte, the writes' included, so each echoed
byte costs twenty-five bytes of script and three bytes of echo about eighty, and the test is then
re-testing the register protocol rather than the echo logic. Worse, it cannot end: the flag that
stops \code{run()} was set by the UART stub, and a transport stub knows nothing of it, so once the
script runs out every \code{STATUS} poll reads zero and \code{run()} spins for ever.

It is also impossible to run \code{EchoNode} over anything \textbf{else} --- the
\code{driver::uart::Stub} of \chapref{6}, a loopback double, a future driver variant, a driver over a
different transport --- because the type is nailed down at compile time.

\textbf{What has to be built first:} the whole test harness of \chapref{8}, one layer too low, and a
transport stub taught to set the stop flag, a double doing an application's job, before a single line
of application logic can be checked.

\textbf{The fix.} One word:

\begin{cppcode}
driver::uart::Interface& myUart;
\end{cppcode}

\textbf{The decision of \chapref{6} it makes concrete.} \textit{(1 mark)} That
\code{driver::uart::Interface} is \textbf{abstract so that the application codes against a promise
rather than against a driver}. The interface exists for exactly this: the application and its tests
are written against the contract, so the same object runs unchanged over the stub in a host test and
over the real \code{Uart} on the bench, and arrives at the bench already proven. Depending on the
concrete class throws the abstraction away while still paying for it --- the vtable, the indirection
and the extra header are all still there, buying nothing.
